\documentclass[12pt, a4paper]{academics}

\academicsCourse{计算机系统基础}{Introduction to Computer Systems}
\academicsReport{浮点数的表示}{2026 年 4 月 27 日}
\academicsArthur{华博文}{19240212}

\begin{document}

\makeacademicscover

\section{实验环境}

\subsection{操作系统}

macOS 26.4 arm64，Darwin 25.4.0，Apple M5。

\subsection{编程工具}

Visual Studio Code 1.118.1，NeoVim v0.12.0，Apple clang 21.0.0。

\subsection{其他工具}

\LaTeX{}。

\section{实验目的}

\begin{enumerate}
    \item 熟悉浮点数值在计算机内部的表示方式；
    \item 掌握浮点数与整型之间的转换和比较现象；
    \item 认识正无穷、负无穷、NaN、\mintinline{c}{+0}、\mintinline{c}{-0} 和非规格化数的位级特征；
    \item 完成 \mintinline{c}{fpower2} 函数的按位实现。
\end{enumerate}

\section{实验内容}

\subsection{浮点数精度实例测试}

\begin{minted}{c}
float data[] = {61.419997f,61.419998f,61.419999f,61.42f,61.420001f};
for (int i = 0; i < 5; i++)
    printf("The number is:%.6f\\n", data[i]);
\end{minted}

运行结果：
\begin{minted}{text}
The number is:61.419998
The number is:61.419998
The number is:61.419998
The number is:61.420000
The number is:61.420002
bits: 0x4279a6e6 0x4279a6e6 0x4279a6e6 0x4279a700 0x4279a70c
\end{minted}

分析与解读：
\begin{enumerate}
    \item 单精度浮点数尾数（含隐含位）约 24 位有效二进制位，因此只能表示有限的二进制分数；十进制小数通常不能精确表示为有限二进制小数，造成上面相近十进制常数被映射到相同或相邻的浮点表示。
    \item 输出差异来自四舍五入与二进制近似，查看二进制（或十六进制机器数）可直观了解差别。
\end{enumerate}

实验过程中，在源代码中依次定义 \mintinline{c}{61.419997f}、\mintinline{c}{61.419998f}、\mintinline{c}{61.419999f}、\mintinline{c}{61.42f}、\mintinline{c}{61.420001f}，并使用 \mintinline{c}{printf("%f", ...)} 输出。观察结果可以发现，前 3 个常量在单精度下几乎落到同一个邻近值上，而后两个值则显示出可见的舍入差异，这与浮点表示的有限精度完全一致。

\subsection{浮点与整型关系表达式永真性判断}

\begin{minted}{c}
int i = (1<<24);
printf("i=%d, (int)(float)i=%d\\n", i, (int)((float)i));
\end{minted}

程序运行给出结论（排除 NaN 与无穷）：
\begin{itemize}
    \item \mintinline{c}{i == (int)(float)i}：不是永真（当 \mintinline{c}{i} 超过 float 的精度范围时会失真）。
示例运行输出（在当前机器上使用默认 64 位编译并运行，输出如下）：
    \item \mintinline{c}{i == (int)(double)i}：对 32 位 \mintinline{c}{int} 在 IEEE-754 双精度下为真（double 精度足以精确表示所有 32 位整数）。
    \item \mintinline{c}{f == (float)(double)f}：在非 NaN 情况下一般为真（\mintinline{c}{float -> double -> float} 保持原值）。
    \item \mintinline{c}{f == -(-f)}：在非 NaN 情况下一般为真（双重取负应恢复原值，包括 \mintinline{c}{+0} 和 \mintinline{c}{-0}）。
\end{itemize}
结合推导可知：第 3 条通常成立，因为双精度的有效尾数足以表示 32 位整数；第 4 条通常成立，因为先提升到 double 再降回 float 不会改变原本能表示的 float 值；第 6 条在非 NaN 情况下也应成立；第 1、2、5、7 条都可能因精度损失、截断或舍入而失败。

\subsection{浮点数特殊值调试实验}

\begin{minted}{c}
float finf = 5e38f; printf("%f\n", finf);
union { unsigned u; float f; } t; t.f = finf;
printf("finf 0x%08x\n", t.u);
\end{minted}

运行输出：
\begin{minted}{text}
inf inf inf
-inf -inf -inf
5.000000 0.100000 
-5.000000 -0.100000
0.00000000000000000000000000000000000000000000000000
-0.00000000000000000000000000000000000000000000000000
0.000000 -0.000000
nan nan
inf 
finf1 0x7f800000
finf2 0x7f800000
finf3 0x7f800000
fninf1 0xff800000
fninf2 0xff800000
fninf3 0xff800000
fzero 0x00000000
fnzero 0x80000000
fnormal1 0x40a00000
fnormal2 0x3dcccccd
fnnormal1 0xc0a00000
fnnormal2 0xbdcccccd
ffrac 0x000116c2
fnfrac 0x800116c2
fnan1 0x7fc00000
fnan2 0xffc00000
finf4 0x7f800000
\end{minted}

对应机理如下：\mintinline{c}{4e38f}、\mintinline{c}{5e38f}、\mintinline{c}{6e38f} 都超过单精度可表示范围，因此都被舍入为正无穷大；相应负数变为负无穷大。\mintinline{c}{0.1f} 和 \mintinline{c}{-0.1f} 只能近似表示，打印到 20 位小数时可以看出尾部误差。\mintinline{c}{1e-40f} 落入非规格化数范围，因此其机器数指数全 0、尾数非 0；\mintinline{c}{finf1 + fninf1} 形成 \mintinline{c}{NaN}，再取负号仍是 \mintinline{c}{NaN}。

\subsection{实验填表}

\begin{center}
\begin{tabular}{|c|c|c|}
\hline
变量 & 机器数（十六进制表示） & 输出值 \\
\hline
finf1 & \mintinline{c}{0x7f800000} & \mintinline{c}{inf} \\
\hline
finf2 & \mintinline{c}{0x7f800000} & \mintinline{c}{inf} \\
\hline
finf3 & \mintinline{c}{0x7f800000} & \mintinline{c}{inf} \\
\hline
fninf1 & \mintinline{c}{0xff800000} & \mintinline{c}{-inf} \\
\hline
fninf2 & \mintinline{c}{0xff800000} & \mintinline{c}{-inf} \\
\hline
fninf3 & \mintinline{c}{0xff800000} & \mintinline{c}{-inf} \\
\hline
fzero & \mintinline{c}{0x00000000} & \mintinline{c}{0.000000} \\
\hline
fnzero & \mintinline{c}{0x80000000} & \mintinline{c}{-0.000000} \\
\hline
fnormal1 & \mintinline{c}{0x40a00000} & \mintinline{c}{5.000000} \\
\hline
fnormal2 & \mintinline{c}{0x3dcccccd} & \mintinline{c}{0.100000} \\
\hline
fnnormal1 & \mintinline{c}{0xc0a00000} & \mintinline{c}{-5.000000} \\
\hline
fnnormal2 & \mintinline{c}{0xbdcccccd} & \mintinline{c}{-0.100000} \\
\hline
ffrac & \mintinline{c}{0x000116c2} & \mintinline{c}{0.000000} \\
\hline
fnfrac & \mintinline{c}{0x800116c2} & \mintinline{c}{-0.000000} \\
\hline
fnan1 & \mintinline{c}{0x7fc00000} & \mintinline{c}{nan} \\
\hline
fnan2 & \mintinline{c}{0xffc00000} & \mintinline{c}{nan} \\
\hline
finf4 & \mintinline{c}{0x7f800000} & \mintinline{c}{inf} \\
\hline
\end{tabular}
\end{center}

表中 \mintinline{c}{ffrac} 和 \mintinline{c}{fnfrac} 的数值依赖编译器与运行平台对非规格化数的处理，但在 IEEE-754 单精度下，\mintinline{c}{1e-40f} 一般会落在非规格化区间，十六进制通常接近 \mintinline{c}{0x000116c2} 一类的小尾数形式。

\subsection{实现 \mintinline{c}{fpower2} 函数}

函数说明：计算 $2^{x}$ 的单精度 \mintinline{c}{float} 表示（按位构造）。

\inputminted{c}{../src/fpower2.c}

\begin{itemize}
    \item 当 $x<-149$ 时下溢，返回 0.0；
    \item 当 $-149\le x \le -127$ 时返回非规格化数，尾数按 $2^{x+149}$ 左移构造；
    \item 当 $-126\le x \le 127$ 时返回规格化数，阶码为 $x+127$，尾数为 0；
    \item 当 $x>127$ 时返回 +∞（阶码全 1，尾数为 0）。
\end{itemize}

结合 IEEE-754 编码规则可知，单精度浮点数由 1 位符号位、8 位阶码和 23 位尾数组成。对于 $2^x$ 而言，数值本身没有尾数，因此规格化结果的尾数部分为 0；当阶码不足以编码时，数值会进入非规格化区间；当指数过大时则直接进入无穷大。实现时采用按位拼接后再通过 \mintinline{c}{u2f} 转换，可以避免依赖运行时的浮点计算过程。

\subsection{实验步骤与说明}

按照实验要求，完成过程可概括为以下几步：

\begin{enumerate}
        \item 编辑源程序，分别完成精度测试、关系表达式测试、特殊值测试和 \mintinline{c}{fpower2} 实现；
        \item 使用 \mintinline{bash}{gcc -g -o test4_3 test4_3.c} 或 \mintinline{bash}{clang -g -o test4_3 test4_3.c} 编译；
        \item 运行程序并截图保存输出结果；
        \item 用 \mintinline{bash}{objdump -S test4_3 > test4_3.txt} 反汇编并分析机器码；
        \item 用 \mintinline{bash}{gdb test4_3} 进行断点调试、查看寄存器和栈帧；
        \item 根据实验环境填写机器数表格与调试记录。
\end{enumerate}

\section{结论}

本实验通过浮点数精度测试、浮点与整型关系判断、特殊值观察和 \mintinline{c}{fpower2} 函数实现，系统验证了 IEEE-754 单精度浮点数的表示规律。实验表明，十进制小数通常无法精确转化为二进制浮点数，浮点与整型混合运算时需要特别关注隐式转换和精度损失；同时，正无穷、负无穷、NaN 和非规格化数都能从机器数位型上直接识别。通过按位构造 \mintinline{c}{fpower2}，进一步加深了对阶码、尾数与边界条件的理解。

\appendix
\section{附录}

\subsection{\mil{precision.c}}
\inputminted{c}{../src/precision.c}

\subsection{\mil{relations.c}}
\inputminted{c}{../src/relations.c}

\subsection{\mil{specials.c}}
\inputminted{c}{../src/specials.c}

\subsection{\mil{fpower2.c}}
\inputminted{c}{../src/fpower2.c}

\end{document}
